% EM structure and teaching revision 20260928
\subsection{非保守电场下的电压表读数}
静电场中，两点电势差与路径无关；变化磁场产生的感生电场满足 $\oint\boldsymbol E\cdot\dd\boldsymbol l=-\dot\Phi_B$，其沿两条不同连接路径的线积分之差一般不为零。
设路径 $\Gamma_1,\Gamma_2$ 都从 $A$ 到 $B$，按 $\Gamma_1$ 去、沿 $\Gamma_2$ 返回组成闭合回路，则
\[
 \int_{\Gamma_1}\boldsymbol E\cdot\dd\boldsymbol l-
 \int_{\Gamma_2}\boldsymbol E\cdot\dd\boldsymbol l=-\frac{\dd\Phi_B}{\dd t}.
\]
因此读电压表时，除了接点 $A,B$，还要给出引线围成的路径以及是否穿过变化磁通区域。把引线绞合或贴近布置，是减小附加回路面积与感应误差的一种方式。
计算可把各段电阻上的电压降和感应电动势按指定回路统一列入方程。应避免给整个空间强行指定一个能代表全部电场的单值静电势；只有感应贡献可忽略或已明确划入集中电动势元件时，普通节点电势法才可直接使用。
测量模型还假定电压表内阻足够大、对原回路电流影响可忽略；若不满足此条件，表支路本身也须参与电路求解。

